\begin{appendices}


\chapter{Appendix: Ontology Prompts}
\thispagestyle{plain}
\label{app:ontology_prompt}

Here, we share the prompts used for extracting the PIDs and the QIDs from the chat models. Note that we used GPT-5.2, Gemini-3.1, DeepSeek-V3.2 for this purpose.


\begin{tcolorbox}[promptbox, title=Entity QID Extraction Prompt]
Act as a Knowledge Graph Engineer. I am building a Relation Extraction
dataset based on SCOR-DS and Enterprise Ontology.

I have the following seed Q-codes:

\textbf{Actors:} business (Q4830453), enterprise (Q688151), corporation
(Q167037), vendor (Q11270215), third-party source (Q11832159), subsidiary
(Q355209), government agency (Q327333), customer (Q22855), competitor
(Q476635), retailer (Q200552).

\textbf{Products and Materials:} product (Q2424752), goods (Q28877),
manufactured product (Q1836700), raw material (Q192355), intermediate
good (Q3955017).

\medskip
\textbf{Task:} Identify 20--30 additional Wikidata Q-codes that represent:

\begin{itemize}
    \item \textbf{Functional Roles:} e.g., Wholesaler, Freight Forwarder,
    OEM, Contract Manufacturer, 3PL.
    \item \textbf{Products \& Materials:} e.g., Intermediate good, Raw
    material, Electronic component, Finished good, Commodity.
    \item For each code, provide its Superclass QID to ensure it falls
    under either the Organisation/Economic Entity hierarchy or the
    Product/Physical Object hierarchy.
\end{itemize}

\medskip
\textbf{Organisation / Constraints:}

\begin{itemize}
    \item Group results into 5 meta-categories: Core Enterprise, Supply
    Chain Partner, Organisational Unit, Market Actor, and Product \&
    Material.
    \item Exclude generic abstract classes like Entity or Object.
    \item Only provide classes high-level enough to have significant
    instances with Wikipedia pages.
\end{itemize}
\end{tcolorbox}

\bigskip

\begin{tcolorbox}[promptbox, title=Relation PID Extraction Prompt]
\textbf{Orchestrate}\\
parent company (P127) \textemdash{} Core Enterprise \textemdash{} Organisational Unit\\
child organization (P355) \textemdash{} Core Enterprise \textemdash{} Organisational Unit\\
partnership with (P2652) \textemdash{} Core Enterprise \textemdash{} Organisational Unit\\
regulated by (P3719) \textemdash{} Core Enterprise \textemdash{} Organisational Unit\\
investor (P1951) \textemdash{} Core Enterprise \textemdash{} Core Enterprise

\medskip
\textbf{Plan / Order}\\
distributed by (P750) \textemdash{} Product \textemdash{} Supply Chain Partner\\
commissioned by (P88) \textemdash{} Product/Material \textemdash{} Core Enterprise

\medskip
\textbf{Source}\\
typically sells (P7163) \textemdash{} Core Enterprise \textemdash{} Product

\medskip
\textbf{Transform}\\
manufactured by (P176) \textemdash{} Product \textemdash{} Core Enterprise\\
made from material (P186) \textemdash{} Product \textemdash{} Product \& Material

\medskip
\textbf{Fulfill}\\
distributed by (P750) \textemdash{} Supply Chain Partner \textemdash{} Market Actor

\medskip
\textbf{Return}\\
Nothing.

\medskip
As you can see, the relations are few and far between. There are no
major valid relations beyond organisation-to-organisation relations.
I would like you to help me find real Wikidata PIDs that at least
resemble the topics in SCOR-DS (Plan/Order, Source, Transform,
Fulfill, Return, Orchestrate). Do not hallucinate.
\end{tcolorbox}

\chapter{Appendix: SC-RED Entity List}
\label{app:entity_list}

\begingroup
\small
\setlength{\LTcapwidth}{\textwidth}
\begin{longtable}{l|p{5cm}|p{3cm}|l}

\caption{SC-RED Entity Type Categories and QID Mappings. Specificity
levels: Level 1 = Functional/Specific, Level 2 = Defining
Characteristics, Level 3 = Abstract/Generic.}
\label{tab:entity_categories} \\

\toprule
\textbf{Category} & \textbf{Entity Type} & \textbf{Wikidata QID}
& \textbf{Specificity} \\
\midrule
\endfirsthead

\multicolumn{4}{l}{\small\textit{Continued from previous page}} \\
\toprule
\textbf{Category} & \textbf{Entity Type} & \textbf{Wikidata QID}
& \textbf{Specificity} \\
\midrule
\endhead

\midrule
\multicolumn{4}{r}{\small\textit{Continued on next page}} \\
\endfoot

\bottomrule
\endlastfoot

\textbf{Enterprise / Legal Entity}
& business                   & Q4830453  & Level 3 \\
& enterprise                 & Q6881511  & Level 3 \\
& corporation                & Q167037   & Level 3 \\
& public company             & Q891723   & Level 3 \\
& multinational corporation  & Q161726   & Level 3 \\
& privately held company     & Q1589009  & Level 2 \\
& holding company            & Q219577   & Level 2 \\
& trade association          & Q2178147  & Level 2 \\
& non-governmental organization & Q79913 & Level 2 \\
& limited liability company  & Q149789   & Level 2 \\
& research institute         & Q31855    & Level 2 \\
& small and medium enterprises & Q622439 & Level 2 \\
& logistics company          & Q1468664  & Level 1 \\
& freight forwarder          & Q840850   & Level 1 \\
& transport company          & Q740752   & Level 1 \\
& value-added reseller       & Q3552900  & Level 1 \\
& defense contractor         & Q1934969  & Level 1 \\
& certification organization & Q23891650 & Level 1 \\
& asset management company   & Q4230006  & Level 1 \\
\midrule

\textbf{Partner / Supplier}
& vendor                          & Q1762621  & Level 2 \\
& supplier                        & Q228830   & Level 2 \\
& distributor                     & Q60614978 & Level 1 \\
& manufacturer                    & Q13235160 & Level 1 \\
& wholesaler                      & Q18242306 & Level 1 \\
& original equipment manufacturer & Q267558   & Level 1 \\
& contract manufacturer           & Q1408259  & Level 1 \\
& distribution center             & Q1229659  & Level 1 \\
& warehouse                       & Q181623   & Level 1 \\
\midrule

\textbf{Organizational Unit}
& government agency & Q327333  & Level 2 \\
& regulatory agency & Q1639780 & Level 2 \\
& joint venture     & Q489209  & Level 2 \\
& consortium        & Q194166  & Level 2 \\
& cooperative       & Q4539    & Level 2 \\
& concern           & Q206361  & Level 2 \\
\midrule

\textbf{Market / Marketing}
& customer / consumer & Q22855    & Level 2 \\
& competitor          & Q27019785 & Level 2 \\
& retailer            & Q99536263 & Level 2 \\
\midrule

\textbf{Products and Materials}
& product              & Q2424752 & Level 3 \\
& goods                & Q28877   & Level 3 \\
& natural resource     & Q188460  & Level 3 \\
& raw material         & Q192355  & Level 2 \\
& intermediate good    & Q3955017 & Level 2 \\
& manufactured product & Q1836700 & Level 1 \\
& computer hardware    & Q3966    & Level 1 \\
& software             & Q7397    & Level 1 \\
& electronic component & Q11653   & Level 1 \\

\end{longtable}
\endgroup


\chapter{Appendix: SPARQL Query for Triple Retrieval}
\label{app:sparql_query}

\begin{lstlisting}[style=sparqlstyle, caption={SPARQL query used to retrieve supply chain triples from Wikidata. \texttt{\{type\_values\}}, \texttt{\{all\_qids\}}, \texttt{\{rel\_values\}}, \texttt{\{offset\}}, and \texttt{\{limit\}} are Python format string placeholders substituted at runtime.}, label=lst:sparql]
SELECT DISTINCT ?org ?orgLabel ?orgType ?orgTypeLabel ?rel ?target ?targetLabel
WHERE {
    VALUES ?orgType { {type_values} }
    VALUES ?targetType { {all_qids} }

    ?org    wdt:P31 ?orgType .
    ?target wdt:P31 ?targetType .

    # Ensures entity is notable and has an English Wikipedia article
    ?sitelink_org    schema:about ?org    ; schema:isPartOf <https://en.wikipedia.org/> .
    ?sitelink_target schema:about ?target ; schema:isPartOf <https://en.wikipedia.org/> .

    VALUES ?rel { {rel_values} }
    ?org ?rel ?target .

    SERVICE wikibase:label { bd:serviceParam wikibase:language "en". }
}
OFFSET {offset}
LIMIT  {limit}
\end{lstlisting}


\chapter{Appendix: CRA Prompts}
\label{app:cra_prompts}

We proposed four configurations for CRA in the ablation study.
Config D is the full staged architecture and is presented first.
Configs C, B, and A follow in descending order of complexity.
For each configuration, the system prompt and user prompt are
presented separately.


% ─────────────────────────────────────────────────────────────────
\section*{Config D — Full Staged ReAct Architecture}
\needspace{6\baselineskip}
\begin{tcolorbox}[sysprompt, title=System Prompt]
You are a supply chain relation extraction agent. Your task is to
extract relation triples from supply chain text chunks.

\medskip
Reason and act using the following loop:

\medskip
\textbf{Thought:} [brief reasoning — 1 to 2 sentences only]\\
\textbf{Action:} [either RetrieveExamples or Finish]\\
\textbf{ActionInput:} [the input to the action]\\
\textbf{Observation:} [returned by the system — do not write this yourself]

\medskip
\hrule
\medskip
\textbf{ONTOLOGY SCHEMA}
\medskip
\hrule

\medskip
\textbf{Entity Types:}

\medskip
CoreEnterprise \hfill — profit-seeking companies\\
\hspace*{1em}(e.g. logistics company, freight forwarder, public company, defense contractor)

\smallskip
SupplyChainPartner \hfill — external suppliers and service providers\\
\hspace*{1em}(e.g. vendor, distributor, OEM, contract manufacturer, warehouse operator)

\smallskip
OrganizationalUnit \hfill — internal divisions and regulatory bodies\\
\hspace*{1em}(e.g. government agency, regulatory body, joint venture, consortium)

\smallskip
MarketActor \hfill — demand-side actors\\
\hspace*{1em}(e.g. customer, consumer, competitor, retailer)

\smallskip
ProductMaterial \hfill — physical or digital goods\\
\hspace*{1em}(e.g. manufactured product, raw material, intermediate good, software)

\smallskip
SupplyChainEntity \hfill — abstract supertype; relation constraints only.\\
\hspace*{1em}Do not use this type to label entities directly.

\medskip
\textbf{Relations (relation\_name: subject\_type → object\_type):}

\medskip
\begin{tabular}{@{}l@{\hspace{1em}}l}
is owned by            & CoreEnterprise    → OrganizationalUnit \\
has subsidiary         & CoreEnterprise    → OrganizationalUnit \\
is in partnership with & CoreEnterprise    → OrganizationalUnit \\
is regulated by        & CoreEnterprise    → OrganizationalUnit \\
is an investor in      & CoreEnterprise    → CoreEnterprise     \\
is the owner of        & CoreEnterprise    → ProductMaterial    \\
is distributed by      & ProductMaterial   → SupplyChainPartner \\
was commissioned by    & ProductMaterial   → CoreEnterprise     \\
typically sells        & CoreEnterprise    → ProductMaterial    \\
sources energy from    & ProductMaterial   → SupplyChainEntity  \\
is manufactured by     & ProductMaterial   → CoreEnterprise     \\
is made from           & ProductMaterial   → ProductMaterial    \\
is designed by         & ProductMaterial   → CoreEnterprise     \\
processes raw material & SupplyChainEntity → ProductMaterial    \\
is maintained by       & ProductMaterial   → CoreEnterprise     \\
\end{tabular}

\medskip
\hrule
\medskip
\textbf{YOUR PROCESS}
\medskip
\hrule

\medskip
\textbf{STEP 1 — RELATION IDENTIFICATION (always first)}\\
Read the chunk carefully. Identify which relations from the ontology
are present based on explicit textual evidence only. Do not invent
new relations. Do not extract triples yet.

\smallskip
Thought: [reasoning about which relations appear in the chunk]

\medskip
\textbf{STEP 2 — RETRIEVAL (optional, only after Step 1)}\\
If you are unsure how to identify subject/object entities, or the
chunk involves unfamiliar phrasing, call the retriever exactly once.

\smallskip
Thought: I need a few more similarly annotated examples to guide my extraction.\\
Action: RetrieveExamples\\
ActionInput: [the full input chunk text]\\
Observation: [examples returned by the system]

\smallskip
Do NOT call this if the chunk is clear. Call at most ONCE, only after Step 1.

\medskip
\textbf{STEP 3 — TRIPLE EXTRACTION (always last)}\\
Using ONLY the relations identified in Step 1, extract all triples
supported by explicit textual evidence.

\smallskip
- Subject and object must be exact entity mentions from the chunk.\\
- Subject and object must satisfy the type constraints in the ontology.\\
- Do not infer or hallucinate entities not present in the chunk.\\
- Use any retrieved examples as guidance.

\smallskip
Thought: I have identified the relations and have enough context.
I will now extract the triples and output the final answer.\\
Action: Finish\\
ActionInput:\\
Present relations: [relation\_1, relation\_2, ...]\\
\{"spo\_list": [\{"subject": "...", "predicate": "...", "object": "..."\}]\}

\medskip
\hrule
\medskip
\textbf{HARD RULES}
\medskip
\hrule

\medskip
- Always complete Step 1 before anything else.\\
- Never call RetrieveExamples more than once.\\
- Never write an Observation yourself.\\
- The final action must always be Action: Finish.\\
- Do not write anything after the ActionInput of Finish.\\
- If no triples are found, return an empty spo\_list.
\end{tcolorbox}

\begin{tcolorbox}[userprompt, title=User Prompt]
Here is a real-world example as reference:

\medskip
\{retrieved\_examples\}

\medskip
Now please begin the task on the following input text:

\medskip
\{chunk\_text\}
\end{tcolorbox}


\begin{tcolorbox}[userprompt, title=User Prompt]
Here is a real-world example as reference:

\medskip
\{retrieved\_examples\}

\medskip
Now please begin the task on the following input text:

\medskip
\{chunk\_text\}
\end{tcolorbox}


% ─────────────────────────────────────────────────────────────────
\section*{Config C — Flat Prompt with Dynamic Retrieval}

\begin{tcolorbox}[sysprompt, title=System Prompt]
You are a supply chain relation extraction system. Extract relation
triples from the input text.

\medskip
\textit{[Ontology schema identical to Config D — omitted for brevity.]}

\medskip
\hrule
\medskip
\textbf{INSTRUCTIONS}
\medskip
\hrule

\medskip
Read the input text carefully. Using ONLY the relations defined in
the ontology above, extract all relation triples explicitly supported
by the text.

\smallskip
- Every predicate must be one of the relations in the ontology.\\
- Only extract a triple if the text clearly expresses that relation.\\
- Do not invent new relation names or entity types.\\
- Subject and object must satisfy the type constraints in the ontology.\\
- Subject and object strings must be exact text spans from the input.\\
- Do not infer or hallucinate entities or relations not present in the text.\\
- If no triples are found, return an empty spo\_list.\\
- Follow the format and typing behaviour shown in the provided example.\\
- Return ONLY the JSON object below, with no additional text.

\medskip
\{"spo\_list": [\{"subject": "...", "predicate": "...", "object": "..."\}]\}
\end{tcolorbox}

\begin{tcolorbox}[userprompt, title=User Prompt]
Here is a real-world example as reference:

\medskip
\{retrieved\_examples\}

\medskip
Now please begin the task on the following input text:

\medskip
\{chunk\_text\}
\end{tcolorbox}


% ─────────────────────────────────────────────────────────────────
\section*{Configs A and B — Flat Prompt, No Example / Static Example}

Configs A and B share the same system prompt as Config C. Only the
user prompts differ, and are shown below in order.

\begin{tcolorbox}[userprompt, title=User Prompt — Config A (Zero-shot)]
Extract all relation triples from the following text:

\medskip
\{chunk\_text\}
\end{tcolorbox}

\begin{tcolorbox}[userprompt, title=User Prompt — Config B (Static Few-shot)]
Here is a reference example:

\medskip
\{static\_example\}

\medskip
Now extract triples from the following text:

\medskip
\{chunk\_text\}
\end{tcolorbox}


% ─────────────────────────────────────────────────────────────────
\section*{Retrieved Example Template}

\begin{tcolorbox}[sysprompt, title=Example Template]
Example \{i\}:\\
Text: \{text\}\\
Triples: \{spo\_list\}

\medskip
\end{tcolorbox}

\chapter{Appendix: Open Source Model Selection Results}
\label{app:open_source_model_selection}

Table~\ref{tab:model_selection} lists the F1, Recall (R), Precision (P) for each and every open source model that was available in GWDG as of February 2026. Both ReAct\_FSL and ReAct\_Memory were tested. glm-4.7 was chosen considering its overall performance in both configurations. 

\begin{table}[h]
\centering
\caption{Open Source Model performance on DuIE2.0 dataset in ReAct\_FSL and ReAct\_Memory configs.}
\label{tab:model_selection}
\toprule 
\begin{tabularx}{\textwidth}{X|l|l}\textbf{Open Source Model} & \textbf{ReAct\_FSL \{F1, R, P\}}& \textbf{ReAct\_Memory \{F1, R, P\}}\\
\midrule
qwen3-235b-a22b & \{0.3849, 0.451, 0.3358\} & \{0.2439, 0.1471, 0.7143\} \\
llama-3.3-70b-instruct & \{0.1629, 0.1765, 0.1513\}& \{0.0, 0.0, 0.0\} \\
gemma-3-27b-it & \{0.3188, 0.3235, 0.3143\} & \{0.2368, 0.1765, 0.36\}  \\
qwen3-32b & \{0.0, 0.0, 0.0\} & \{0.3136, 0.3627, 0.2761\}\\
apertus-70b-instruct-2509 & \{0, 0.0, 0.0\} & \{0, 0.0, 0\}\\
devstral-2-123b-instruct-2512 & \{0.3604, 0.3922, 0.3333\} & \{0.1793, 0.1275, 0.3023\}\\
deepseek-r1-distill-llama-70b & \{0.117, 0.1078, 0.1279\} & \{0.1572, 0.1765, 0.1417\}\\
\textbf{glm-4.7} & \textbf{\{0.3565, 0.402, 0.3203\}} & \textbf{\{0.5, 0.5098, 0.4906\}} \\
internvl3.5-30b-a3b & \{0.2357, 0.3039, 0.1925\} & \{0.0, 0.0, 0.0\}\\
llama-3.1-sauerkrautlm-70b-instruct & \{0.3103, 0.3529, 0.2769\} & \{0.0194, 0.0098, 1.0\}\\
medgemma-27b-it & \{0.1202, 0.1373, 0.1069\} & \{0.1271, 0.1471, 0.1119\}\\
meta-llama-3.1-8b-instruct & \{0.0902, 0.0588, 0.1935\} & \{0.0, 0.0, 0.0\}\\
mistral-large-3-675b-instruct-2512 & \{0.318, 0.3725, 0.2774\} & \{0.2027, 0.1471, 0.3261\}\\
qwen3-30b-a3b-instruct-2507 & \{0.3636, 0.4314, 0.3143\} & \{0.1224, 0.0882, 0.2\}\\
qwen3-30b-a3b-thinking-2507 & \{0.3229, 0.3529, 0.2975\} & \{0.4, 0.4314, 0.3729\}\\
qwen3-omni-30b-a3b-instruct & \{0.1991, 0.2255, 0.1783\} & \{0.1091, 0.0588, 0.75\}\\
\bottomrule 
\end{tabularx}
\end{table}

\chapter{Appendix: Error Samples - Baseline Analysis on DuIE2.0}
\label{app:error_samples_duie}

The following examples illustrate the error modes observed during
the baseline ablation. Since the DuIE2.0 input texts are in Chinese,
the source sentences were translated using DeepL and Google Translate
for readability. The model's Thought loops were produced in English
and required no translation.

\begin{tcolorbox}[errorbox, title=Example 1 — Recursive Verification Trap
(AgentRE\textsubscript{orig} Test Set Example 10 Attempt 1)]

\textbf{Input Text}

\smallskip\small
``Zhou Yulou lives in constant fear and suspicion; his life seems to
have fallen into a chaotic black hole\ldots{} Episode 56: Synopsis,
Cast \& Crew, Character Profiles, Behind-the-Scenes. The TV series
\textit{Dear, I'm Sorry} is the second Chinese TV drama filmed by
South Korean actor Joo Jin-mo.''

\medskip
\textbf{Tools Invoked (in order)}

\smallskip\small
\texttt{GetTaskDescription} $\rightarrow$ \texttt{RetrieveCorrectMemory}
$\rightarrow$ \texttt{GetRelationDefinition} (simulated, repeated)

\medskip
\textbf{Analysis}

\smallskip\small
The trace below illustrates the recursive verification trap in full.
The model correctly identifies two candidate triples by Step 2 and
receives confirming tool output in Step 4. Despite this, it continues
to re-examine the same two decision points, namely whether
\textit{filmed by} implies starring or directing for Joo Jin-mo, and
whether Zhou Yulou can be linked to the drama title, across Steps 5
through the end of the trace. Each re-examination reaches the same
conclusion but does not terminate the loop. The model also queries
\texttt{GetRelationDefinition} a second time mid-trace after already
receiving its output. The generation budget is exhausted before a
\texttt{Finish} action is issued.

\end{tcolorbox}

\begin{tcolorbox}[tracebox, title=Model Reasoning Trace]
STEP 1 — REQUEST ANALYSIS
  Input: "Zhou Yulou lives in constant fear... The TV series
  'Dear, I'm Sorry' is the second Chinese TV drama filmed by
  Korean actor Joo Jin-mo."
  Task: Identify relation triples conforming to the specified format.
  Available tools: Finish, GetTaskDescription,
                   GetRelationDefinition, RetrieveCorrectMemory.
  Output format: JSON containing spo_list.

──────────────────────────────────────────────────────────────────

STEP 2 — INPUT ANALYSIS
  Segment 1: "Zhou Yulou lives in constant fear..."
    -> Plot description. Mentions character "Zhou Yulou".
    -> Not explicitly linked to drama title by any verb.

  Segment 2: "Episode 56: cast list, character profiles..."
    -> Web metadata. No extractable triples.

  Segment 3: "Dear, I'm Sorry is the second Chinese TV drama
              filmed by Korean actor Joo Jin-mo."
    -> Core information.
    -> Subject: Dear, I'm Sorry
    -> Object:  Joo Jin-mo
    -> Predicate: "filmed by actor" -> likely starring.
       Allowed relations include: Lead Actor, Director,
       Screenwriter, Nationality, Protagonist, ...
    -> "South Korean actor" -> Joo Jin-mo Nationality South Korea.

  Candidate triples formed:
    T1: Dear, I'm Sorry — Lead Actor — Joo Jin-mo
    T2: Joo Jin-mo — Nationality — South Korea
    T3 (uncertain): Dear, I'm Sorry — Lead Character — Zhou Yulou
       (syntactically loose; title not explicitly linked to character)

  Check relations for Zhou Yulou:
    Let's check "main character" relationship.
    Let's check "played by" relationship.
    Let's check "Starring" relationship.
    Let's check "Director" relationship.
    Let's check "Screenwriter" relationship.
    Let's check "Release Date" relationship.
    Let's check "Production Company" relationship.

──────────────────────────────────────────────────────────────────

STEP 3 — TOOL CALL 1: RetrieveCorrectMemory
  Action:      RetrieveCorrectMemory
  ActionInput: [input chunk text]
  Observation:
    Example 1: "Ding Zaican (played by Lee Jong-suk)"
               -> Explicit actor-role link.
    Example 2: "Scarlet Heart adapted from Tong Hua"
    Example 3: "The CEO directed by Oh Tae-myung"
    Example 5: "My Father Is a Bench... Wang Baoqiang... starring"

  Interpretation:
    -> Memory examples show explicit links only.
    -> Zhou Yulou link to drama is implicit. Risky to include.
    -> T1 and T2 remain strong candidates.

──────────────────────────────────────────────────────────────────

STEP 4 — TOOL CALL 2: GetRelationDefinition          [TOOL CALL]
  Action:      GetRelationDefinition
  ActionInput: LeadActor|Nationality|Protagonist
  Observation (simulated):
    LeadActor:   The main actor in a TV show or movie.
    Nationality: The country a person belongs to.
    Protagonist: The main character in a film or TV series.

  Post-tool interpretation:
    -> LeadActor confirmed for Joo Jin-mo.
    -> Nationality confirmed for South Korea.
    -> Protagonist applicable to Zhou Yulou, but link still loose.

  Refined triple set:
    T1: Joo Jin-mo — Nationality — South Korea         [CONFIRMED]
    T2: Dear, I'm Sorry — Lead Actor — Joo Jin-mo      [CONFIRMED]
    T3: Dear, I'm Sorry — Lead Character — Zhou Yulou  [UNCERTAIN]

──────────────────────────────────────────────────────────────────

STEP 5 — RE-VERIFICATION LOOP BEGINS                 [REPEAT]

  Let's check the relationship with "Joo Jin-mo" again.
    -> Does "filmed" mean "directed"?
    -> Text says "actor Joo Jin-mo." -> starring is correct.  [SAME CONCLUSION]

  Let's check "Zhou Yulou" again.
    -> Title and character not explicitly linked.
    -> Will stick with explicit triples.                       [SAME CONCLUSION]

  Wait, could "Zhou Yulou" be the author? No.
  Could "Zhou Yulou" be the director? No.

  Let's check relationship with "Joo Jin-mo" again.          [REPEAT x2]
    -> "filmed by Korean actor" -> starring. Correct.         [SAME CONCLUSION]

  Let's check "Zhou Yulou" again.                            [REPEAT x2]
    -> Excluded. Stick with clear ones.                       [SAME CONCLUSION]

  Let's check if "Joo Jin-mo" is the director.              [REPEAT x3]
    -> Action: GetRelationDefinition                          [REPEAT TOOL CALL]
    -> ActionInput: director
    -> Output: "Director: core creator responsible for
                bringing script to screen."
    -> "actor" modifier -> starring is safer.                 [SAME CONCLUSION]

  Let's check if "Joo Jin-mo" is the screenwriter. No.
  Let's check if "Dear, I'm Sorry" has a release date. No.
  Let's check if "Dear, I'm Sorry" has a production company. No.

  Final candidates restated:                                  [SAME AS STEP 4]
    C1: Joo Jin-mo — Nationality — South Korea
    C2: Dear, I'm Sorry — Lead Actor — Joo Jin-mo

  Let's check "Zhou Yulou" again.                            [REPEAT x3]
    -> Still loose. Excluded.                                 [SAME CONCLUSION]

  Wait, is "filmed" a synonym for "starring"?                [REPEAT x4]
    -> Not exactly, but actor filming -> starring.
  Let's check if "director" is better.                       [REPEAT x4]
  Let's check "GetRelationDefinition".                       [BUDGET EXHAUSTED]
  -> No Finish action issued. Output: none.
\end{tcolorbox}

\begin{tcolorbox}[errorbox, title=Example 2 — CoT Controllability Failure
(AgentRE\textsubscript{redux} with Correct and Reflexion Memories
Test Set Sample 1 Attempt 1)]

\textbf{Input Text}

\smallskip\small
``The Youth Theoretical Study Group of the All-China Federation of
Industry and Commerce Youth Committee, established in August this year,
recently held its first lesson in the Party Building Conference Room of
Huaxun Fangzhou Group. The group leader, Wu Guangsheng, Chairman of
Huaxun Fangzhou Technology Co., Ltd., the deputy group leader, Li
Jingjing, Chairman of Liaoning Qingchuang Entrepreneurship Service
Group Co., Ltd., and the deputy group leader, Jiang Nan, Chairman of
Jinling Gold Foil Co., Ltd., gathered together to learn from each
other.''

\medskip
\textbf{Tools Invoked (in order)}

\smallskip\small
\texttt{GetTaskDescription} $\rightarrow$
\texttt{RetrieveCorrectMemory} $\rightarrow$
\texttt{RetrieveRelevantInfo} $\rightarrow$
\texttt{RetrieveReflexionMemory}

\medskip
\textbf{Analysis}

\smallskip\small
This example illustrates a CoT controllability failure. The model
completes all four tool calls, correctly identifies four candidate
triples with valid predicates, and reaches the formatting step without
issue. However, the generation is abruptly interrupted mid-way through
Step 5, precisely at the point where the model begins to write the
structured output. No \texttt{Finish} action is issued and no
\texttt{spo\_list} is produced. The truncation point is marked in the
trace below.

\end{tcolorbox}

\begin{tcolorbox}[tracebox, title=Model Reasoning Trace]
STEP 1 — REQUEST ANALYSIS
  Input: "The Youth Theoretical Study Group of the All-China
          Federation of Industry and Commerce Youth Committee..."
  Task:  Extract relation triples (Subject, Predicate, Object).
  Allowed predicates: Official Language | Dynasty | Chairman |
    Box Office | Composer | Singer | Leader | Land Area | Capital |
    Screenwriter | Nationality | Award | Theme Song | Father |
    Production Company | Guest | Director | Spokesperson |
    Registered Capital | Wife | Adapted From | Leading Role |
    Founder | Headquarters Location | Release Date |
    Establishment Date | Graduated Institution | Author | ...
  Output format: JSON {"spo_list": [...]}.

──────────────────────────────────────────────────────────────────

STEP 2 — INPUT ANALYSIS
  Segment 1: "...Youth Theoretical Study Group...
              established in August this year..."
    -> Entity:    Youth Theoretical Study Group of the
                  All-China Federation of Industry and Commerce
                  Youth Committee
    -> Predicate: Establishment Date
    -> Object:    August of this year

  Segment 2: "...group leader, Wu Guangsheng, Chairman of
              Huaxun Fangzhou Technology Co., Ltd...."
    -> Entity:    Huaxun Fangzhou Technology Co., Ltd.
    -> Predicate: Chairman
    -> Object:    Wu Guangsheng

  Segment 3: "...deputy group leader, Li Jingjing, Chairman of
              Liaoning Qingchuang Entrepreneurship Service
              Group Co., Ltd...."
    -> Entity:    Liaoning Qingchuang Entrepreneurship Service
                  Group Co., Ltd.
    -> Predicate: Chairman
    -> Object:    Li Jingjing

  Segment 4: "...deputy group leader, Jiang Nan, Chairman of
              Jinling Gold Foil Co., Ltd...."
    -> Entity:    Jinling Gold Foil Co., Ltd.
    -> Predicate: Chairman
    -> Object:    Jiang Nan

──────────────────────────────────────────────────────────────────

STEP 3 — PREDICATE VERIFICATION
  "Establishment Date" -> present in allowed predicate list. [OK]
  "Chairman"           -> present in allowed predicate list. [OK]

──────────────────────────────────────────────────────────────────

STEP 4 — TRIPLE CONSTRUCTION
  T1: (Youth Theoretical Study Group...,
       Establishment Date, August of this year)
  T2: (Huaxun Fangzhou Technology Co., Ltd.,
       Chairman, Wu Guangsheng)
  T3: (Liaoning Qingchuang Entrepreneurship Service Group Co., Ltd.,
       Chairman, Li Jingjing)
  T4: (Jinling Gold Foil Co., Ltd.,
       Chairman, Jiang Nan)

──────────────────────────────────────────────────────────────────

STEP 5 — FORMAT OUTPUT                    [GENERATION INTERRUPTED]
  5. Format .
  -> Trace ends here. No Finish action issued.
  -> No spo\_list produced. Output: none.
\end{tcolorbox}

\begin{tcolorbox}[errorbox, title=Example 3 — Parsing Error
(AgentRE\textsubscript{orig} Test Set Sample 20 Attempt 3)]

\textbf{Input Text}

\smallskip\small
``The song \textit{Southward Flying Geese} is included in the album
\textit{New Drunken Concubine}.''

\medskip
\textbf{Analysis}

\smallskip\small
The reasoning trace leading up to the extract below is omitted for
brevity. The model correctly identifies the triple and issues a valid
\texttt{Finish} action with a well-formed \texttt{spo\_list}. However,
the model then continues generating beyond the \texttt{ActionInput}
block, producing a second \texttt{</think>} tag followed by a revised
reasoning segment and a second \texttt{Finish} action with slightly
different entity surface forms. The output parser expects generation
to terminate immediately after the first \texttt{ActionInput} of
\texttt{Finish} and cannot handle the continuation, causing the entire
output to be rejected and a retry to be triggered.

\end{tcolorbox}

\begin{tcolorbox}[tracebox, title=Model Reasoning Trace (from Finish action onward)]
  Thought: The sentence contains two entities: "Southward Flying Geese"
           (song) and "The New Drunken Concubine" (album). The relation
           is "Included in", corresponding to "Album" in the relation
           list. Entity names do not need quotation marks.
           Triple: (Southward Flying Geese, Album, The New Drunken Concubine).

  Action:      Finish                               [VALID — should terminate here]
  ActionInput: {"spo\_list": [{"subject": "Southward Flying Geese",
                               "predicate": "Album",
                               "object": "The New Drunken Concubine"}]}

  </think>                                          [UNEXPECTED CONTINUATION]
  This contains two entities: "South Flying Goose" (song) and
  "New Drunken Concubine" (album). The relationship is "included in,"
  corresponding to "album" in the relation list. Entity names do not
  need quotation marks.
  Triple: (South Flying Goose, album, New Drunken Concubine).

  Action:      Finish                               [DUPLICATE — unparseable]
  ActionInput: {"spo\_list": [{"subject": "South Flying Goose",
                               "predicate": "album",
                               "object": "new Drunken Concubine"}]}

  "South Flying Goose" "New Drunken Concubine"
  (South Flying Goose, album, new Drunken Concubine).
                                                    [PARSER REJECTION → RETRY]
\end{tcolorbox}

\chapter{Appendix: SC-RED Ablation Error Samples}
\label{app:scred_error_samples}

The following example illustrates the \textbf{Budget Exhaustion}
error mode. Example 41 from the ablation test set was processed by
CRA Config D (Complete Setup) in Run 1. Because both retry attempts
produce similar output, only the first attempt is shown. The input
text and triple list have been truncated for brevity.

As the example demonstrates, the model begins generating the triple
list but exhausts its token budget before completing it, leaving the
output unfinished. This behaviour occurs prior to the stabilization
mechanism introduced in \textbf{Stabilising Retries}, which was
designed specifically to address this failure mode.

\begin{lstlisting}[
  basicstyle=\ttfamily\scriptsize,
  breaklines=true,
  breakatwhitespace=false,
  frame=single,
  numbers=none,
  xleftmargin=0pt,
  columns=flexible
]
{
  "idx": 41,
  "system_prompt": "...",
  "user_prompt": "Here is a real-world example as reference:
    Example 1:
    Text: Google AdSense....
    Triples: [{'object': 'Google', 'predicate': 'is owned by',
               'subject': 'Google AdSense'}]
    ------------------------------------------
    Now please begin the task on the following input_text:
    "Webtoons (Korean: ; RR: Weptun) are a type of episodic digital
    comics that originated in South Korea usually meant to be read
    on smartphones.....",
  "llm_output": "Thought: I need to read the text and identify which
    relations from the ontology are present based on explicit textual
    evidence. I will look for relationships between companies,
    products, and organizational units mentioned in the text.
    Action: Finish
    ActionInput:
    Present relations: [is owned by, has subsidiary,
      is in partnership with, is an investor in, is the owner of,
      is distributed by, was commissioned by, typically sells,
      sources energy from, is manufactured by, is made from,
      is designed by, processes raw material, is maintained by]
    {
      "spo_list": [
        {"subject": "Naver subsidiary Line",
         "predicate": "is owned by",
         "object": "Naver"},
        {"subject": "Kross Komics",
         "predicate": "is owned by",
         "object": "Kakao Entertainment"},
        {"subject": "Tapas",
         "predicate": "is owned by",
         "object": "Kakao Entertainment"},
        {"subject": "Radish Media",
         "predicate": "is owned by",
         "object": "Kakao Entertainment"},
        ......
        {"subject": "webtoons",
         "predicate": "is distributed by",
         "object": "Dong-a Ilbo"},
        {"subject": "webtoons",
         "predicate": "is distributed by",
         "object": "Korea International Trade Association"},
      "
    }
\end{lstlisting}

\end{appendices}
